\documentclass[
  doc,
  floatsintext,
  longtable,
  a4paper,
  nolmodern,
  notxfonts,
  notimes,
  12pt,
  colorlinks=true,linkcolor=blue,citecolor=blue,urlcolor=blue]{apa7}

\usepackage{amsmath}
\usepackage{amssymb}

\geometry{inner=1in, outer=1in}
\fancyhfoffset[LE,RO]{0cm}


\usepackage[bidi=default]{babel}
\babelprovide[main,import]{spanish}


% get rid of language-specific shorthands (see #6817):
\let\LanguageShortHands\languageshorthands
\def\languageshorthands#1{}

\RequirePackage{longtable}
\RequirePackage{threeparttablex}

\makeatletter
\renewcommand{\paragraph}{\@startsection{paragraph}{4}{\parindent}%
	{0\baselineskip \@plus 0.2ex \@minus 0.2ex}%
	{-.5em}%
	{\normalfont\normalsize\bfseries\typesectitle}}

\renewcommand{\subparagraph}[1]{\@startsection{subparagraph}{5}{0.5em}%
	{0\baselineskip \@plus 0.2ex \@minus 0.2ex}%
	{-\z@\relax}%
	{\normalfont\normalsize\bfseries\itshape\hspace{\parindent}{#1}\textit{\addperi}}{\relax}}
\makeatother




\usepackage{longtable, booktabs, multirow, multicol, colortbl, hhline, caption, array, float, xpatch}
\usepackage{subcaption}


\renewcommand\thesubfigure{\Alph{subfigure}}
\setcounter{topnumber}{2}
\setcounter{bottomnumber}{2}
\setcounter{totalnumber}{4}
\renewcommand{\topfraction}{0.85}
\renewcommand{\bottomfraction}{0.85}
\renewcommand{\textfraction}{0.15}
\renewcommand{\floatpagefraction}{0.7}

\usepackage{tcolorbox}
\tcbuselibrary{listings,theorems, breakable, skins}
\usepackage{fontawesome5}

\definecolor{quarto-callout-color}{HTML}{909090}
\definecolor{quarto-callout-note-color}{HTML}{0758E5}
\definecolor{quarto-callout-important-color}{HTML}{CC1914}
\definecolor{quarto-callout-warning-color}{HTML}{EB9113}
\definecolor{quarto-callout-tip-color}{HTML}{00A047}
\definecolor{quarto-callout-caution-color}{HTML}{FC5300}
\definecolor{quarto-callout-color-frame}{HTML}{ACACAC}
\definecolor{quarto-callout-note-color-frame}{HTML}{4582EC}
\definecolor{quarto-callout-important-color-frame}{HTML}{D9534F}
\definecolor{quarto-callout-warning-color-frame}{HTML}{F0AD4E}
\definecolor{quarto-callout-tip-color-frame}{HTML}{02B875}
\definecolor{quarto-callout-caution-color-frame}{HTML}{FD7E14}

%\newlength\Oldarrayrulewidth
%\newlength\Oldtabcolsep


\usepackage{hyperref}



\usepackage{color}
\usepackage{fancyvrb}
\newcommand{\VerbBar}{|}
\newcommand{\VERB}{\Verb[commandchars=\\\{\}]}
\DefineVerbatimEnvironment{Highlighting}{Verbatim}{commandchars=\\\{\}}
% Add ',fontsize=\small' for more characters per line
\usepackage{framed}
\definecolor{shadecolor}{RGB}{241,243,245}
\newenvironment{Shaded}{\begin{snugshade}}{\end{snugshade}}
\newcommand{\AlertTok}[1]{\textcolor[rgb]{0.68,0.00,0.00}{#1}}
\newcommand{\AnnotationTok}[1]{\textcolor[rgb]{0.37,0.37,0.37}{#1}}
\newcommand{\AttributeTok}[1]{\textcolor[rgb]{0.40,0.45,0.13}{#1}}
\newcommand{\BaseNTok}[1]{\textcolor[rgb]{0.68,0.00,0.00}{#1}}
\newcommand{\BuiltInTok}[1]{\textcolor[rgb]{0.00,0.23,0.31}{#1}}
\newcommand{\CharTok}[1]{\textcolor[rgb]{0.13,0.47,0.30}{#1}}
\newcommand{\CommentTok}[1]{\textcolor[rgb]{0.37,0.37,0.37}{#1}}
\newcommand{\CommentVarTok}[1]{\textcolor[rgb]{0.37,0.37,0.37}{\textit{#1}}}
\newcommand{\ConstantTok}[1]{\textcolor[rgb]{0.56,0.35,0.01}{#1}}
\newcommand{\ControlFlowTok}[1]{\textcolor[rgb]{0.00,0.23,0.31}{\textbf{#1}}}
\newcommand{\DataTypeTok}[1]{\textcolor[rgb]{0.68,0.00,0.00}{#1}}
\newcommand{\DecValTok}[1]{\textcolor[rgb]{0.68,0.00,0.00}{#1}}
\newcommand{\DocumentationTok}[1]{\textcolor[rgb]{0.37,0.37,0.37}{\textit{#1}}}
\newcommand{\ErrorTok}[1]{\textcolor[rgb]{0.68,0.00,0.00}{#1}}
\newcommand{\ExtensionTok}[1]{\textcolor[rgb]{0.00,0.23,0.31}{#1}}
\newcommand{\FloatTok}[1]{\textcolor[rgb]{0.68,0.00,0.00}{#1}}
\newcommand{\FunctionTok}[1]{\textcolor[rgb]{0.28,0.35,0.67}{#1}}
\newcommand{\ImportTok}[1]{\textcolor[rgb]{0.00,0.46,0.62}{#1}}
\newcommand{\InformationTok}[1]{\textcolor[rgb]{0.37,0.37,0.37}{#1}}
\newcommand{\KeywordTok}[1]{\textcolor[rgb]{0.00,0.23,0.31}{\textbf{#1}}}
\newcommand{\NormalTok}[1]{\textcolor[rgb]{0.00,0.23,0.31}{#1}}
\newcommand{\OperatorTok}[1]{\textcolor[rgb]{0.37,0.37,0.37}{#1}}
\newcommand{\OtherTok}[1]{\textcolor[rgb]{0.00,0.23,0.31}{#1}}
\newcommand{\PreprocessorTok}[1]{\textcolor[rgb]{0.68,0.00,0.00}{#1}}
\newcommand{\RegionMarkerTok}[1]{\textcolor[rgb]{0.00,0.23,0.31}{#1}}
\newcommand{\SpecialCharTok}[1]{\textcolor[rgb]{0.37,0.37,0.37}{#1}}
\newcommand{\SpecialStringTok}[1]{\textcolor[rgb]{0.13,0.47,0.30}{#1}}
\newcommand{\StringTok}[1]{\textcolor[rgb]{0.13,0.47,0.30}{#1}}
\newcommand{\VariableTok}[1]{\textcolor[rgb]{0.07,0.07,0.07}{#1}}
\newcommand{\VerbatimStringTok}[1]{\textcolor[rgb]{0.13,0.47,0.30}{#1}}
\newcommand{\WarningTok}[1]{\textcolor[rgb]{0.37,0.37,0.37}{\textit{#1}}}

\providecommand{\tightlist}{%
  \setlength{\itemsep}{0pt}\setlength{\parskip}{0pt}}
\usepackage{longtable,booktabs,array}
\newcounter{none} % for unnumbered tables
\usepackage{calc} % for calculating minipage widths
% Correct order of tables after \paragraph or \subparagraph
\usepackage{etoolbox}
\makeatletter
\patchcmd\longtable{\par}{\if@noskipsec\mbox{}\fi\par}{}{}
\makeatother
% Allow footnotes in longtable head/foot
\IfFileExists{footnotehyper.sty}{\usepackage{footnotehyper}}{\usepackage{footnote}}
\makesavenoteenv{longtable}

\usepackage{graphicx}
\makeatletter
\newsavebox\pandoc@box
\newcommand*\pandocbounded[1]{% scales image to fit in text height/width
  \sbox\pandoc@box{#1}%
  \Gscale@div\@tempa{\textheight}{\dimexpr\ht\pandoc@box+\dp\pandoc@box\relax}%
  \Gscale@div\@tempb{\linewidth}{\wd\pandoc@box}%
  \ifdim\@tempb\p@<\@tempa\p@\let\@tempa\@tempb\fi% select the smaller of both
  \ifdim\@tempa\p@<\p@\scalebox{\@tempa}{\usebox\pandoc@box}%
  \else\usebox{\pandoc@box}%
  \fi%
}
% Set default figure placement to htbp
\def\fps@figure{htbp}
\makeatother







\usepackage{newtx}

\defaultfontfeatures{Scale=MatchLowercase}
\defaultfontfeatures[\rmfamily]{Ligatures=TeX,Scale=1}





\title{Cadenas de texto}


\shorttitle{CADENAS DE TEXTO}


\usepackage{etoolbox}



\ccoppy{\textcopyright~2026}



\author{Edison Achalma}



\affiliation{
{Escuela Profesional de Economía, Universidad Nacional de San Cristóbal
de Huamanga}}




\leftheader{Achalma}

\date{2026-10-20}


\abstract{Este abstract será actualizado una vez que se complete el
contenido final del artículo. }



\authornote{\par{\addORCIDlink{Edison Achalma}{0000-0001-6996-3364}} 
\par{ }
\par{   El autor no tiene conflictos de interés que revelar.    Los
roles de autor se clasificaron utilizando la taxonomía de roles de
colaborador (CRediT; https://credit.niso.org/) de la siguiente
manera:  Edison Achalma:   conceptualización, metodología, análisis
formal, investigación, recursos, curación de
datos, redacción, visualización, supervisión, administración del
proyecto}
\par{La correspondencia relativa a este artículo debe dirigirse a Edison
Achalma, Escuela Profesional de Economía, Universidad Nacional de San
Cristóbal de Huamanga, Portal Independencia N°
57, Ayacucho, AYA 5001, Perú, Email: \href{mailto:elmer.achalma.09@unsch.edu.pe}{elmer.achalma.09@unsch.edu.pe}}
}

\makeatletter
\let\endoldlt\endlongtable
\def\endlongtable{
\hline
\endoldlt
}
\makeatother

\urlstyle{same}



\hypersetup{
  pdftitle={Cadenas de texto},
  pdfauthor={Autor}
}
\makeatletter
\@ifpackageloaded{caption}{}{\usepackage{caption}}
\AtBeginDocument{%
\ifdefined\contentsname
  \renewcommand*\contentsname{Tabla de contenidos}
\else
  \newcommand\contentsname{Tabla de contenidos}
\fi
\ifdefined\listfigurename
  \renewcommand*\listfigurename{Lista de Figuras}
\else
  \newcommand\listfigurename{Lista de Figuras}
\fi
\ifdefined\listtablename
  \renewcommand*\listtablename{Lista de Tablas}
\else
  \newcommand\listtablename{Lista de Tablas}
\fi
\ifdefined\figurename
  \renewcommand*\figurename{Figura}
\else
  \newcommand\figurename{Figura}
\fi
\ifdefined\tablename
  \renewcommand*\tablename{Tabla}
\else
  \newcommand\tablename{Tabla}
\fi
}
\@ifpackageloaded{float}{}{\usepackage{float}}
\floatstyle{ruled}
\@ifundefined{c@chapter}{\newfloat{codelisting}{h}{lop}}{\newfloat{codelisting}{h}{lop}[chapter]}
\floatname{codelisting}{Listado}
\newcommand*\listoflistings{\listof{codelisting}{Lista de Listados}}
\makeatother
\makeatletter
\makeatother
\makeatletter
\@ifpackageloaded{caption}{}{\usepackage{caption}}
\@ifpackageloaded{subcaption}{}{\usepackage{subcaption}}
\makeatother
\makeatletter
\@ifpackageloaded{fontawesome5}{}{\usepackage{fontawesome5}}
\makeatother

% From https://tex.stackexchange.com/a/645996/211326
%%% apa7 doesn't want to add appendix section titles in the toc
%%% let's make it do it
\makeatletter
\xpatchcmd{\appendix}
  {\par}
  {\addcontentsline{toc}{section}{\@currentlabelname}\par}
  {}{}
\makeatother

%% Disable longtable counter
%% https://tex.stackexchange.com/a/248395/211326

\usepackage{etoolbox}

\makeatletter
\patchcmd{\LT@caption}
  {\bgroup}
  {\bgroup\global\LTpatch@captiontrue}
  {}{}
\patchcmd{\longtable}
  {\par}
  {\par\global\LTpatch@captionfalse}
  {}{}
\apptocmd{\endlongtable}
  {\ifLTpatch@caption\else\addtocounter{table}{-1}\fi}
  {}{}
\newif\ifLTpatch@caption
\makeatother

\begin{document}

\maketitle


\hypertarget{toc}{}
\tableofcontents
\newpage
\section[Introduction]{Cadenas de texto}

\setcounter{secnumdepth}{3}

\setlength\LTleft{0pt}




\subsection{🎯 Motivación}\label{motivaciuxf3n}

Las cadenas de texto (strings) son uno de los tipos de datos más
utilizados en programación. Desde mostrar mensajes al usuario hasta
procesar datos de archivos o bases de datos, el manejo de texto es
fundamental. Python ofrece una forma potente y flexible de trabajar con
strings, con una sintaxis clara y numerosos métodos para manipularlos.

\subsection{💡 Intuición}\label{intuiciuxf3n}

Imagina un string como una \textbf{fila de caracteres} donde cada
posición tiene un número (índice). Es como un tren donde cada vagón es
una letra, y puedes acceder a cualquier vagón por su número. Python
trata los strings como secuencias inmutables: una vez creados, no puedes
cambiar los caracteres individuales, pero puedes crear nuevos strings a
partir de ellos.

\subsection{🧠 Modelo mental}\label{modelo-mental}

\begin{itemize}
\tightlist
\item
  \textbf{Strings son secuencias}: cada carácter tiene una posición
  (índice), empezando desde 0.
\item
  \textbf{Inmutables}: no puedes cambiar un carácter directamente; creas
  uno nuevo.
\item
  \textbf{Indexación}: accedes con \texttt{{[}i{]}} (i desde 0). Índices
  negativos empiezan desde el final.
\item
  \textbf{Multilínea}: con triple comilla (\texttt{"""} o
  \texttt{\textquotesingle{}\textquotesingle{}\textquotesingle{}})
  preserva saltos de línea.
\item
  \textbf{Inmutabilidad}: los métodos devuelven nuevos strings, no
  modifican el original.
\end{itemize}

\subsection{📐 Desarrollo}\label{desarrollo}

\subsubsection{1. Creación de Strings}\label{creaciuxf3n-de-strings}

\begin{Shaded}
\begin{Highlighting}[]
\CommentTok{\# Comillas simples o dobles}
\NormalTok{nombre }\OperatorTok{=} \StringTok{"Ana"}
\NormalTok{apellido }\OperatorTok{=} \StringTok{\textquotesingle{}García\textquotesingle{}}

\CommentTok{\# Strings con comillas internas}
\NormalTok{texto1 }\OperatorTok{=} \StringTok{"It\textquotesingle{}s a beautiful day"}
\NormalTok{texto2 }\OperatorTok{=} \StringTok{\textquotesingle{}Ella dijo "Hola"\textquotesingle{}}

\CommentTok{\# Strings multilínea (con triple comilla)}
\NormalTok{poema }\OperatorTok{=} \StringTok{"""Roses are red,}
\StringTok{Violets are blue,}
\StringTok{Python is awesome,}
\StringTok{And so are you!"""}
\BuiltInTok{print}\NormalTok{(poema)}
\end{Highlighting}
\end{Shaded}

\subsubsection{2. Acceso a caracteres
(indexación)}\label{acceso-a-caracteres-indexaciuxf3n}

\begin{Shaded}
\begin{Highlighting}[]
\NormalTok{texto }\OperatorTok{=} \StringTok{"Python"}
\BuiltInTok{print}\NormalTok{(texto[}\DecValTok{0}\NormalTok{])  }\CommentTok{\# \textquotesingle{}P\textquotesingle{}}
\BuiltInTok{print}\NormalTok{(texto[}\DecValTok{1}\NormalTok{])  }\CommentTok{\# \textquotesingle{}y\textquotesingle{}}
\BuiltInTok{print}\NormalTok{(texto[}\DecValTok{2}\NormalTok{])  }\CommentTok{\# \textquotesingle{}t\textquotesingle{}}

\CommentTok{\# Índices negativos (desde el final)}
\BuiltInTok{print}\NormalTok{(texto[}\OperatorTok{{-}}\DecValTok{1}\NormalTok{])  }\CommentTok{\# \textquotesingle{}n\textquotesingle{}}
\BuiltInTok{print}\NormalTok{(texto[}\OperatorTok{{-}}\DecValTok{2}\NormalTok{])  }\CommentTok{\# \textquotesingle{}o\textquotesingle{}}

\CommentTok{\# ❌ Error: índice fuera de rango}
\CommentTok{\# print(texto[10])  \# IndexError}
\end{Highlighting}
\end{Shaded}

\subsubsection{3. Longitud de un string}\label{longitud-de-un-string}

\begin{Shaded}
\begin{Highlighting}[]
\NormalTok{texto }\OperatorTok{=} \StringTok{"Hola Mundo"}
\BuiltInTok{print}\NormalTok{(}\BuiltInTok{len}\NormalTok{(texto))  }\CommentTok{\# 10 (incluye el espacio)}
\end{Highlighting}
\end{Shaded}

\subsubsection{4. Slicing (rebanado)}\label{slicing-rebanado}

Extrae partes de un string:

\begin{Shaded}
\begin{Highlighting}[]
\NormalTok{texto }\OperatorTok{=} \StringTok{"Python Programming"}

\CommentTok{\# texto[inicio:fin]  (fin no incluido)}
\BuiltInTok{print}\NormalTok{(texto[}\DecValTok{0}\NormalTok{:}\DecValTok{6}\NormalTok{])   }\CommentTok{\# "Python"}
\BuiltInTok{print}\NormalTok{(texto[}\DecValTok{7}\NormalTok{:}\DecValTok{18}\NormalTok{])  }\CommentTok{\# "Programming"}

\CommentTok{\# Desde el inicio}
\BuiltInTok{print}\NormalTok{(texto[:}\DecValTok{6}\NormalTok{])    }\CommentTok{\# "Python"}

\CommentTok{\# Hasta el final}
\BuiltInTok{print}\NormalTok{(texto[}\DecValTok{7}\NormalTok{:])    }\CommentTok{\# "Programming"}

\CommentTok{\# Con índices negativos}
\BuiltInTok{print}\NormalTok{(texto[}\OperatorTok{{-}}\DecValTok{11}\NormalTok{:])  }\CommentTok{\# "Programming"}
\BuiltInTok{print}\NormalTok{(texto[}\OperatorTok{{-}}\DecValTok{11}\NormalTok{:}\OperatorTok{{-}}\DecValTok{1}\NormalTok{]) }\CommentTok{\# "Programmin"}
\end{Highlighting}
\end{Shaded}

\subsubsection{5. Iteración sobre
strings}\label{iteraciuxf3n-sobre-strings}

\begin{Shaded}
\begin{Highlighting}[]
\NormalTok{palabra }\OperatorTok{=} \StringTok{"Python"}
\ControlFlowTok{for}\NormalTok{ letra }\KeywordTok{in}\NormalTok{ palabra:}
    \BuiltInTok{print}\NormalTok{(letra)}
\CommentTok{\# Salida: P, y, t, h, o, n}

\CommentTok{\# Con índice y carácter}
\ControlFlowTok{for}\NormalTok{ i, letra }\KeywordTok{in} \BuiltInTok{enumerate}\NormalTok{(palabra):}
    \BuiltInTok{print}\NormalTok{(}\SpecialStringTok{f"Posición }\SpecialCharTok{\{}\NormalTok{i}\SpecialCharTok{\}}\SpecialStringTok{: }\SpecialCharTok{\{}\NormalTok{letra}\SpecialCharTok{\}}\SpecialStringTok{"}\NormalTok{)}
\end{Highlighting}
\end{Shaded}

\subsubsection{\texorpdfstring{6. Verificar contenido con
\texttt{in}}{6. Verificar contenido con in}}\label{verificar-contenido-con-in}

\begin{Shaded}
\begin{Highlighting}[]
\NormalTok{frase }\OperatorTok{=} \StringTok{"El gato está en el tejado"}

\CommentTok{\# Verificar si contiene una subcadena}
\BuiltInTok{print}\NormalTok{(}\StringTok{"gato"} \KeywordTok{in}\NormalTok{ frase)      }\CommentTok{\# True}
\BuiltInTok{print}\NormalTok{(}\StringTok{"perro"} \KeywordTok{in}\NormalTok{ frase)     }\CommentTok{\# False}
\BuiltInTok{print}\NormalTok{(}\StringTok{"tejado"} \KeywordTok{not} \KeywordTok{in}\NormalTok{ frase) }\CommentTok{\# False}

\CommentTok{\# Uso en condicional}
\ControlFlowTok{if} \StringTok{"gato"} \KeywordTok{in}\NormalTok{ frase:}
    \BuiltInTok{print}\NormalTok{(}\StringTok{"¡Hay un gato!"}\NormalTok{)}
\end{Highlighting}
\end{Shaded}

\subsubsection{7. Operaciones básicas con
strings}\label{operaciones-buxe1sicas-con-strings}

\begin{Shaded}
\begin{Highlighting}[]
\CommentTok{\# Concatenación}
\NormalTok{nombre }\OperatorTok{=} \StringTok{"Ana"}
\NormalTok{apellido }\OperatorTok{=} \StringTok{"García"}
\NormalTok{nombre\_completo }\OperatorTok{=}\NormalTok{ nombre }\OperatorTok{+} \StringTok{" "} \OperatorTok{+}\NormalTok{ apellido}
\BuiltInTok{print}\NormalTok{(nombre\_completo)  }\CommentTok{\# "Ana García"}

\CommentTok{\# Repetición}
\NormalTok{eco }\OperatorTok{=} \StringTok{"Hola! "} \OperatorTok{*} \DecValTok{3}
\BuiltInTok{print}\NormalTok{(eco)  }\CommentTok{\# "Hola! Hola! Hola! "}

\CommentTok{\# Comparación (lexicográfica)}
\BuiltInTok{print}\NormalTok{(}\StringTok{"abc"} \OperatorTok{\textless{}} \StringTok{"abd"}\NormalTok{)   }\CommentTok{\# True (c \textless{} d)}
\BuiltInTok{print}\NormalTok{(}\StringTok{"Ana"} \OperatorTok{==} \StringTok{"ana"}\NormalTok{)  }\CommentTok{\# False (diferentes mayúsculas)}
\BuiltInTok{print}\NormalTok{(}\StringTok{"Ana"}\NormalTok{.lower() }\OperatorTok{==} \StringTok{"ana"}\NormalTok{.lower())  }\CommentTok{\# True}
\end{Highlighting}
\end{Shaded}

\subsection{🔬 Ejemplos}\label{ejemplos}

\textbf{Mínimo ---} Creación y acceso:

\begin{Shaded}
\begin{Highlighting}[]
\NormalTok{mensaje }\OperatorTok{=} \StringTok{"Hola Mundo"}
\BuiltInTok{print}\NormalTok{(mensaje[}\DecValTok{0}\NormalTok{])     }\CommentTok{\# H}
\BuiltInTok{print}\NormalTok{(mensaje[}\DecValTok{5}\NormalTok{])     }\CommentTok{\# M}
\BuiltInTok{print}\NormalTok{(}\BuiltInTok{len}\NormalTok{(mensaje))   }\CommentTok{\# 10}
\end{Highlighting}
\end{Shaded}

\textbf{Progresivo ---} Manipulación básica:

\begin{Shaded}
\begin{Highlighting}[]
\NormalTok{nombre }\OperatorTok{=} \StringTok{"Carlos"}
\BuiltInTok{print}\NormalTok{(nombre.upper())  }\CommentTok{\# CARLOS}
\BuiltInTok{print}\NormalTok{(nombre.lower())  }\CommentTok{\# carlos}
\BuiltInTok{print}\NormalTok{(nombre[}\DecValTok{0}\NormalTok{])       }\CommentTok{\# C}
\BuiltInTok{print}\NormalTok{(nombre[}\OperatorTok{{-}}\DecValTok{1}\NormalTok{])      }\CommentTok{\# s}
\BuiltInTok{print}\NormalTok{(}\BuiltInTok{len}\NormalTok{(nombre))     }\CommentTok{\# 6}
\end{Highlighting}
\end{Shaded}

\textbf{Real ---} Sistema de validación de usuario:

\begin{Shaded}
\begin{Highlighting}[]
\CommentTok{\# Datos de usuario}
\NormalTok{nombre\_usuario }\OperatorTok{=} \StringTok{"  Ana\_Perez  "}
\NormalTok{email }\OperatorTok{=} \StringTok{"ANA@GMAIL.COM"}

\CommentTok{\# Limpieza de datos}
\NormalTok{nombre\_limpio }\OperatorTok{=}\NormalTok{ nombre\_usuario.strip()  }\CommentTok{\# "Ana\_Perez"}
\NormalTok{email\_limpio }\OperatorTok{=}\NormalTok{ email.lower()            }\CommentTok{\# "ana@gmail.com"}

\CommentTok{\# Validación}
\ControlFlowTok{if} \StringTok{"@"} \KeywordTok{in}\NormalTok{ email\_limpio:}
    \BuiltInTok{print}\NormalTok{(}\SpecialStringTok{f"Email válido: }\SpecialCharTok{\{}\NormalTok{email\_limpio}\SpecialCharTok{\}}\SpecialStringTok{"}\NormalTok{)}
\ControlFlowTok{else}\NormalTok{:}
    \BuiltInTok{print}\NormalTok{(}\StringTok{"Email inválido"}\NormalTok{)}

\CommentTok{\# Extracción de nombre de usuario}
\ControlFlowTok{if} \StringTok{"\_"} \KeywordTok{in}\NormalTok{ nombre\_limpio:}
\NormalTok{    partes }\OperatorTok{=}\NormalTok{ nombre\_limpio.split(}\StringTok{"\_"}\NormalTok{)}
    \BuiltInTok{print}\NormalTok{(}\SpecialStringTok{f"Nombre: }\SpecialCharTok{\{}\NormalTok{partes[}\DecValTok{0}\NormalTok{]}\SpecialCharTok{\}}\SpecialStringTok{"}\NormalTok{)  }\CommentTok{\# "Ana"}
    \BuiltInTok{print}\NormalTok{(}\SpecialStringTok{f"Apellido: }\SpecialCharTok{\{}\NormalTok{partes[}\DecValTok{1}\NormalTok{]}\SpecialCharTok{\}}\SpecialStringTok{"}\NormalTok{)  }\CommentTok{\# "Perez"}
\end{Highlighting}
\end{Shaded}

\subsection{⚠️ Errores comunes}\label{errores-comunes}

{\def\LTcaptype{none} % do not increment counter
\begin{longtable}[]{@{}
  >{\raggedright\arraybackslash}p{(\linewidth - 4\tabcolsep) * \real{0.2519}}
  >{\raggedright\arraybackslash}p{(\linewidth - 4\tabcolsep) * \real{0.4427}}
  >{\raggedright\arraybackslash}p{(\linewidth - 4\tabcolsep) * \real{0.3053}}@{}}
\toprule\noalign{}
\begin{minipage}[b]{\linewidth}\raggedright
Error
\end{minipage} & \begin{minipage}[b]{\linewidth}\raggedright
Síntoma
\end{minipage} & \begin{minipage}[b]{\linewidth}\raggedright
Solución
\end{minipage} \\
\midrule\noalign{}
\endhead
\bottomrule\noalign{}
\endlastfoot
\textbf{Índice fuera de rango} &
\texttt{IndexError:\ string\ index\ out\ of\ range} & Verifica que el
índice sea \textless{} len(string) \\
\textbf{Confundir \texttt{=} con \texttt{==}} & Asignación en lugar de
comparación & Usa \texttt{==} para comparar \\
\textbf{Modificar string directamente} &
\texttt{TypeError:\ \textquotesingle{}str\textquotesingle{}\ object\ does\ not\ support\ item\ assignment}
& Crea un nuevo string o usa listas \\
\textbf{Olvidar convertir tipos} &
\texttt{TypeError:\ can\ only\ concatenate\ str\ to\ str} & Convierte
con \texttt{str()} o usa f-strings \\
\textbf{Espacios invisibles} & Comparaciones fallan & Usa
\texttt{.strip()} para limpiar \\
\end{longtable}
}

\subsection{✅ Buenas prácticas}\label{buenas-pruxe1cticas}

\begin{enumerate}
\def\labelenumi{\arabic{enumi}.}
\item
  \textbf{Usa f-strings para formateo}:

\begin{Shaded}
\begin{Highlighting}[]
\NormalTok{nombre }\OperatorTok{=} \StringTok{"Ana"}
\NormalTok{edad }\OperatorTok{=} \DecValTok{25}
\NormalTok{mensaje }\OperatorTok{=} \SpecialStringTok{f"Hola, }\SpecialCharTok{\{}\NormalTok{nombre}\SpecialCharTok{\}}\SpecialStringTok{. Tienes }\SpecialCharTok{\{}\NormalTok{edad}\SpecialCharTok{\}}\SpecialStringTok{ años"}
\end{Highlighting}
\end{Shaded}
\item
  \textbf{Limpia strings} antes de comparar:

\begin{Shaded}
\begin{Highlighting}[]
\NormalTok{usuario }\OperatorTok{=} \StringTok{"  ana  "}
\ControlFlowTok{if}\NormalTok{ usuario.strip().lower() }\OperatorTok{==} \StringTok{"ana"}\NormalTok{:}
    \BuiltInTok{print}\NormalTok{(}\StringTok{"Usuario válido"}\NormalTok{)}
\end{Highlighting}
\end{Shaded}
\item
  \textbf{Usa métodos de string} en lugar de bucles manuales:

\begin{Shaded}
\begin{Highlighting}[]
\CommentTok{\# ❌ Menos eficiente}
\NormalTok{resultado }\OperatorTok{=} \StringTok{""}
\ControlFlowTok{for}\NormalTok{ char }\KeywordTok{in} \StringTok{"abc"}\NormalTok{:}
\NormalTok{    resultado }\OperatorTok{+=}\NormalTok{ char.upper()}

\CommentTok{\# ✅ Mejor}
\NormalTok{resultado }\OperatorTok{=} \StringTok{"abc"}\NormalTok{.upper()}
\end{Highlighting}
\end{Shaded}
\item
  \textbf{Prefiere \texttt{in} para búsquedas}:

\begin{Shaded}
\begin{Highlighting}[]
\CommentTok{\# ✅ Claro y eficiente}
\ControlFlowTok{if} \StringTok{"error"} \KeywordTok{in}\NormalTok{ mensaje:}
    \BuiltInTok{print}\NormalTok{(}\StringTok{"Hay un error"}\NormalTok{)}
\end{Highlighting}
\end{Shaded}
\item
  \textbf{Usa slicing} para extraer partes:

\begin{Shaded}
\begin{Highlighting}[]
\NormalTok{email }\OperatorTok{=} \StringTok{"usuario@email.com"}
\NormalTok{nombre }\OperatorTok{=}\NormalTok{ email[:email.index(}\StringTok{"@"}\NormalTok{)]}
\NormalTok{dominio }\OperatorTok{=}\NormalTok{ email[email.index(}\StringTok{"@"}\NormalTok{) }\OperatorTok{+} \DecValTok{1}\NormalTok{:]}
\end{Highlighting}
\end{Shaded}
\end{enumerate}

\subsection{📝 Ejercicio}\label{ejercicio}

\textbf{Intermedio} --- Procesamiento de texto

Escribe un programa que:

\begin{enumerate}
\def\labelenumi{\arabic{enumi}.}
\tightlist
\item
  Tenga un string \texttt{texto} con ``El rápido zorro marrón salta
  sobre el perro perezoso''
\item
  Extraiga y muestre:

  \begin{itemize}
  \tightlist
  \item
    El primer carácter
  \item
    Los últimos 5 caracteres
  \item
    La palabra ``zorro'' (usando slicing)
  \item
    La longitud total
  \item
    El texto en mayúsculas
  \item
    El texto en minúsculas
  \end{itemize}
\item
  Verifique si contiene ``perro'' y ``gato''
\item
  Reemplace ``perro'' por ``gato'' en el texto
\item
  Divida el texto en palabras y muestre cuántas hay
\end{enumerate}

\begin{center}\rule{0.5\linewidth}{0.5pt}\end{center}

\textbf{Solución (no mires hasta haberlo intentado):}

\begin{Shaded}
\begin{Highlighting}[]
\NormalTok{texto }\OperatorTok{=} \StringTok{"El rápido zorro marrón salta sobre el perro perezoso"}

\CommentTok{\# 1. Caracteres}
\BuiltInTok{print}\NormalTok{(}\SpecialStringTok{f"Primer carácter: }\SpecialCharTok{\{}\NormalTok{texto[}\DecValTok{0}\NormalTok{]}\SpecialCharTok{\}}\SpecialStringTok{"}\NormalTok{)}
\BuiltInTok{print}\NormalTok{(}\SpecialStringTok{f"Últimos 5: }\SpecialCharTok{\{}\NormalTok{texto[}\OperatorTok{{-}}\DecValTok{5}\NormalTok{:]}\SpecialCharTok{\}}\SpecialStringTok{"}\NormalTok{)}
\BuiltInTok{print}\NormalTok{(}\SpecialStringTok{f"Palabra \textquotesingle{}zorro\textquotesingle{}: }\SpecialCharTok{\{}\NormalTok{texto[}\DecValTok{10}\NormalTok{:}\DecValTok{15}\NormalTok{]}\SpecialCharTok{\}}\SpecialStringTok{"}\NormalTok{)  }\CommentTok{\# Asumiendo posición fija}

\CommentTok{\# 2. Longitud y caso}
\BuiltInTok{print}\NormalTok{(}\SpecialStringTok{f"Longitud: }\SpecialCharTok{\{}\BuiltInTok{len}\NormalTok{(texto)}\SpecialCharTok{\}}\SpecialStringTok{"}\NormalTok{)}
\BuiltInTok{print}\NormalTok{(}\SpecialStringTok{f"Mayúsculas: }\SpecialCharTok{\{}\NormalTok{texto}\SpecialCharTok{.}\NormalTok{upper()}\SpecialCharTok{\}}\SpecialStringTok{"}\NormalTok{)}
\BuiltInTok{print}\NormalTok{(}\SpecialStringTok{f"Minúsculas: }\SpecialCharTok{\{}\NormalTok{texto}\SpecialCharTok{.}\NormalTok{lower()}\SpecialCharTok{\}}\SpecialStringTok{"}\NormalTok{)}

\CommentTok{\# 3. Búsquedas}
\BuiltInTok{print}\NormalTok{(}\SpecialStringTok{f"¿Contiene \textquotesingle{}perro\textquotesingle{}? }\SpecialCharTok{\{}\StringTok{\textquotesingle{}perro\textquotesingle{}} \KeywordTok{in}\NormalTok{ texto}\SpecialCharTok{\}}\SpecialStringTok{"}\NormalTok{)}
\BuiltInTok{print}\NormalTok{(}\SpecialStringTok{f"¿Contiene \textquotesingle{}gato\textquotesingle{}? }\SpecialCharTok{\{}\StringTok{\textquotesingle{}gato\textquotesingle{}} \KeywordTok{in}\NormalTok{ texto}\SpecialCharTok{\}}\SpecialStringTok{"}\NormalTok{)}

\CommentTok{\# 4. Reemplazo}
\NormalTok{texto\_modificado }\OperatorTok{=}\NormalTok{ texto.replace(}\StringTok{"perro"}\NormalTok{, }\StringTok{"gato"}\NormalTok{)}
\BuiltInTok{print}\NormalTok{(}\SpecialStringTok{f"Texto modificado: }\SpecialCharTok{\{}\NormalTok{texto\_modificado}\SpecialCharTok{\}}\SpecialStringTok{"}\NormalTok{)}

\CommentTok{\# 5. División en palabras}
\NormalTok{palabras }\OperatorTok{=}\NormalTok{ texto.split()}
\BuiltInTok{print}\NormalTok{(}\SpecialStringTok{f"Número de palabras: }\SpecialCharTok{\{}\BuiltInTok{len}\NormalTok{(palabras)}\SpecialCharTok{\}}\SpecialStringTok{"}\NormalTok{)}
\BuiltInTok{print}\NormalTok{(}\SpecialStringTok{f"Primeras 3 palabras: }\SpecialCharTok{\{}\NormalTok{palabras[:}\DecValTok{3}\NormalTok{]}\SpecialCharTok{\}}\SpecialStringTok{"}\NormalTok{)}
\end{Highlighting}
\end{Shaded}

\subsection{🗂️ Síntesis}\label{suxedntesis}

\begin{itemize}
\tightlist
\item
  \textbf{Strings}: secuencias de caracteres, inmutables, indexación
  desde 0.
\item
  \textbf{Comillas}: simples (\texttt{\textquotesingle{}}), dobles
  (\texttt{"}), triples (\texttt{"""} o
  \texttt{\textquotesingle{}\textquotesingle{}\textquotesingle{}}).
\item
  \textbf{Indexación}: \texttt{{[}i{]}} (positivo) y \texttt{{[}-i{]}}
  (negativo desde el final).
\item
  \textbf{Slicing}: \texttt{{[}inicio:fin:paso{]}} para extraer
  subcadenas.
\item
  \textbf{\texttt{len()}}: obtiene la longitud del string.
\item
  \textbf{Métodos comunes}: \texttt{.upper()}, \texttt{.lower()},
  \texttt{.strip()}, \texttt{.split()}, \texttt{.replace()},
  \texttt{.join()}.
\item
  \textbf{\texttt{in}}: verifica si una subcadena está presente.
\item
  \textbf{F-strings}: la forma recomendada de formatear strings con
  variables.
\end{itemize}

\subsection{🔗 Qué sigue}\label{quuxe9-sigue}

\begin{itemize}
\tightlist
\item
  \textbf{Relacionado:} métodos de string, f-strings, expresiones
  regulares
\item
  \textbf{Siguiente tema:} Métodos de strings --- profundiza en las
  herramientas para manipular texto.
\end{itemize}

\section{Publicaciones Similares}\label{publicaciones-similares}

Si te interesó este artículo, te recomendamos que explores otros blogs y
recursos relacionados que pueden ampliar tus conocimientos. Aquí te dejo
algunas sugerencias:

\begin{enumerate}
\def\labelenumi{\arabic{enumi}.}
\tightlist
\item
  \href{https://numerus-scriptum.netlify.app/python/2020-06-19-instalacion-de-anaconda/index.pdf}{\faIcon{file-pdf}}
  \href{https://numerus-scriptum.netlify.app/python/2020-06-19-instalacion-de-anaconda}{Instalacion
  De Anaconda}
\item
  \href{https://numerus-scriptum.netlify.app/python/2020-06-20-configurar-entorno-virtual-python-anaconda/index.pdf}{\faIcon{file-pdf}}
  \href{https://numerus-scriptum.netlify.app/python/2020-06-20-configurar-entorno-virtual-python-anaconda}{Configurar
  Entorno Virtual Python Anaconda}
\item
  \href{https://numerus-scriptum.netlify.app/python/2021-04-17-01-introducion-a-la-programacion-con-python/index.pdf}{\faIcon{file-pdf}}
  \href{https://numerus-scriptum.netlify.app/python/2021-04-17-01-introducion-a-la-programacion-con-python}{01
  Introducion A La Programacion Con Python}
\item
  \href{https://numerus-scriptum.netlify.app/python/2021-05-31-02-variables-expresiones-y-statements-con-python/index.pdf}{\faIcon{file-pdf}}
  \href{https://numerus-scriptum.netlify.app/python/2021-05-31-02-variables-expresiones-y-statements-con-python}{02
  Variables Expresiones Y Statements Con Python}
\item
  \href{https://numerus-scriptum.netlify.app/python/2021-06-07-03-objetos-de-python/index.pdf}{\faIcon{file-pdf}}
  \href{https://numerus-scriptum.netlify.app/python/2021-06-07-03-objetos-de-python}{03
  Objetos De Python}
\item
  \href{https://numerus-scriptum.netlify.app/python/2021-06-14-04-ejecucion-condicional-con-python/index.pdf}{\faIcon{file-pdf}}
  \href{https://numerus-scriptum.netlify.app/python/2021-06-14-04-ejecucion-condicional-con-python}{04
  Ejecucion Condicional Con Python}
\item
  \href{https://numerus-scriptum.netlify.app/python/2021-06-21-05-iteraciones-con-python/index.pdf}{\faIcon{file-pdf}}
  \href{https://numerus-scriptum.netlify.app/python/2021-06-21-05-iteraciones-con-python}{05
  Iteraciones Con Python}
\item
  \href{https://numerus-scriptum.netlify.app/python/2021-08-16-06-funciones-con-python/index.pdf}{\faIcon{file-pdf}}
  \href{https://numerus-scriptum.netlify.app/python/2021-08-16-06-funciones-con-python}{06
  Funciones Con Python}
\item
  \href{https://numerus-scriptum.netlify.app/python/2021-08-23-07-dataframes-con-python/index.pdf}{\faIcon{file-pdf}}
  \href{https://numerus-scriptum.netlify.app/python/2021-08-23-07-dataframes-con-python}{07
  Dataframes Con Python}
\item
  \href{https://numerus-scriptum.netlify.app/python/2021-11-29-08-prediccion-y-metrica-de-performance-con-python/index.pdf}{\faIcon{file-pdf}}
  \href{https://numerus-scriptum.netlify.app/python/2021-11-29-08-prediccion-y-metrica-de-performance-con-python}{08
  Prediccion Y Metrica De Performance Con Python}
\item
  \href{https://numerus-scriptum.netlify.app/python/2021-12-06-09-metodos-de-machine-learning-para-clasificacion-con-python/index.pdf}{\faIcon{file-pdf}}
  \href{https://numerus-scriptum.netlify.app/python/2021-12-06-09-metodos-de-machine-learning-para-clasificacion-con-python}{09
  Metodos De Machine Learning Para Clasificacion Con Python}
\item
  \href{https://numerus-scriptum.netlify.app/python/2021-12-13-10-metodos-de-machine-learning-para-regresion-con-python/index.pdf}{\faIcon{file-pdf}}
  \href{https://numerus-scriptum.netlify.app/python/2021-12-13-10-metodos-de-machine-learning-para-regresion-con-python}{10
  Metodos De Machine Learning Para Regresion Con Python}
\item
  \href{https://numerus-scriptum.netlify.app/python/2025-05-10-visualizacion-de-datos-con-python/index.pdf}{\faIcon{file-pdf}}
  \href{https://numerus-scriptum.netlify.app/python/2025-05-10-visualizacion-de-datos-con-python}{Visualizacion
  De Datos Con Python}
\end{enumerate}

Esperamos que encuentres estas publicaciones igualmente interesantes y
útiles. ¡Disfruta de la lectura!






\end{document}
